\documentclass[aps,prd,10pt,superscriptaddress,nofootinbib,showkeys,amsmath,amssymb]{revtex4-2}

\usepackage{graphicx}
\usepackage{bm}
\usepackage{booktabs}
\usepackage{array}
\usepackage{amsthm}
\usepackage{enumerate}
\usepackage{hyperref}
\usepackage{lineno}

\hypersetup{
  colorlinks=true,
  linkcolor=blue,
  citecolor=blue,
  urlcolor=blue
}

\newtheorem{theorem}{Theorem}
\newtheorem{corollary}{Corollary}

\begin{document}

\title{Lepton Masses from Z\textsubscript{3}-Symmetric Vortex Modes and the Origin of the Koide Formula}

\author{Steve Bico Mujjabi, MD}
\email{msbico@gmail.com}
\homepage{https://github.com/bampita-bico/CDFD-Part-I-Release}
\thanks{Independent researcher; ORCID: \href{https://orcid.org/0009-0001-0556-5516}{0009-0001-0556-5516}.}
\affiliation{VuraLabs, Kampala, Uganda}

\date[]{July 2026}

\begin{abstract}
The Koide formula $Q = (m_e + m_\mu + m_\tau)/(\sqrt{m_e}+\sqrt{m_\mu}+\sqrt{m_\tau})^2
= 2/3$ has been satisfied by the measured lepton masses to better than one part in
$10^4$ since its discovery in 1981, yet no standard-model explanation exists for this
coincidence.
In this paper, the second in the Constraint-Driven Field Theory (CDFT) series,
we show that the Koide formula follows naturally from a specific
$\mathbb{Z}_3$-symmetric mass parametrisation.
The vortex ring established in Paper~I as the CDFT model of the electron
has a unique stable azimuthal excitation at mode number $n = 3$: the trefoil.
We prove as a mathematical theorem that any three quantities constructed from
the $\mathbb{Z}_3$-symmetric Brannen form
$m_k = M(1 + \sqrt{2}\cos(\theta + 2\pi k/3))^2$ ($k = 0,1,2$)
satisfy $Q = 2/3$ exactly, for \emph{all} $M > 0$ and \emph{all} $\theta$.
The proof requires only two elementary trigonometric identities and contains
no free parameters.
We further argue, within the vortex-ring model, that $n = 3$ is the only stable azimuthal mode of a constrained
vortex ring: $n = 1$ is a zero mode, $n = 2$ is unstable, and $n \geq 4$ are
energetically suppressed.
Fitting $M$ and $\theta$ to any two of the three lepton masses, we predict all
three to within $0.015\%$ of PDG values.
The phase $\theta$ is not a free parameter: with Koide automatic, two lepton
masses completely determine $M$ and $\theta$.
Its derivation from CDFT vacuum properties without using measured masses as input
is the central task of Paper~III.

\medskip
\noindent In the extended notation used across the series, the vacuum limit is written
as $\Psi_s = (\Phi/C)\cdot S\cdot M_s$. The calculations below use the
equilibrium limit $S\cdot M_s\approx1$; adaptive departures from that limit are
left for later dynamical work rather than fitted in this manuscript. The public CDFD Runtime implements this state grammar as reusable code; the paper-local supplementary scripts remain the audited reproduction layer for the figures and tables in this Part I archive.
\end{abstract}

\keywords{Constraint-Driven Field Theory, Constraint-Driven Flux Dynamics, adaptive operating ratio, vacuum structure, Mujjabi tests, reproducible physics, Koide formula, lepton masses, Z3 symmetry}

\maketitle
\linenumbers

\section*{Project Context}
This manuscript constitutes Part I (Paper 02) of the multi-part series
\textit{Constraint-Driven Flux Dynamics (CDFD)}. The underlying mathematical
architecture builds directly upon the concepts established in Paper 01 of this
series (Steve Bico Mujjabi, \textit{``CDFD Part I: Fundamental Physics ---
Paper 01: Vortex Stability in a Constrained Transport Vacuum and the Origin of
the Fine-Structure Constant''}, manuscript; paper-specific DOI pending).

\section*{Symbol Conventions}

\noindent\textbf{CDFD physical-vacuum symbol convention.}
\begin{center}
{\small
\begin{tabular}{p{0.16\linewidth}p{0.33\linewidth}p{0.41\linewidth}}
\hline
Symbol & Part I meaning & Physical reading \\
\hline
$J$ or $\Phi$ & Driving flux or throughput & Transport current, field throughput, circulation drive, or generic flux; a subscript fixes the exact channel \\
$C$ & Active constraint or capacity burden & Vacuum transport capacity, geometric confinement, stiffness, boundary resistance, or load-bearing limit \\
$S$ & Surface responsiveness & Local ability of the vacuum or modeled surface to reconfigure under drive \\
$M_s$ & Structural memory & Retained geometry, phase history, stiffness history, or accumulated constraint state \\
$\Psi_s$ & Adaptive operating ratio $(J/C)\cdot S\cdot M_s$ or $(\Phi/C)\cdot S\cdot M_s$ & Local bookkeeping gauge for constrained operation; it is not a new universal constant \\
$R,a,\chi$ & Ring radius, core radius, and aspect ratio $R/a$ & Geometry used to connect vortex structure to dimensionless electromagnetic scales \\
$\rho_0,c,\kappa$ & Vacuum density scale, propagation speed, and transport stiffness & Medium parameters used in stability, pressure, and self-consistency calculations \\
\hline
\end{tabular}
}
\end{center}

\noindent\textbf{Greek-symbol convention.}
\begin{center}
{\small
\begin{tabular}{p{0.16\linewidth}p{0.33\linewidth}p{0.41\linewidth}}
\hline
Symbol & Local role & Interpretation rule \\
\hline
$\alpha$ & Fine-structure constant unless a local equation states otherwise & In the Part I physics papers, un-subscripted $\alpha$ means the electromagnetic fine-structure constant \\
$\beta,\gamma$ & Local response, stiffness, coupling, or correction coefficients & Equation-local coefficients; subscripts bind them to the named mode, channel, or approximation \\
$\theta,\phi,\Phi_c$ & Phase, orientation, or chirality angles & Used for torus-knot orientation, vacuum phase, or chiral phase geometry as defined in the nearby equation \\
$\delta,\epsilon$ & Perturbation, residual, tolerance, or small-offset scales & Local stability margins, numerical residuals, or experimental tolerances \\
$\lambda,\mu,\nu$ & Rate, mass, mode, or frequency labels & Meaning is fixed by the equation, subscript, and physics context where the symbol appears \\
$\rho,\sigma,\tau$ & Density, spread/noise, or time-scale terms & Local density, variance, stress, relaxation, or transport-memory quantities \\
$\Omega,\Lambda$ & State/domain space or integrated measure where defined & Reserved for explicitly defined state spaces, spectra, or synthesis-level quantities \\
\hline
\end{tabular}
}
\end{center}


\section{Introduction}

In 1981, Yoshio Koide noticed that the masses of the charged leptons satisfy
a remarkably tight relation~\cite{Koide1983}:
\begin{equation}
  Q \;=\; \frac{m_e + m_\mu + m_\tau}
               {(\sqrt{m_e} + \sqrt{m_\mu} + \sqrt{m_\tau})^2}
  \;=\; \frac{2}{3}\,.
  \label{eq:koide}
\end{equation}
Using current PDG values~\cite{PDG2022}, $Q = 0.666661$, which agrees with $2/3$
to one part in $10^4$.
More strikingly, when Koide first wrote the formula, only $m_e$ and $m_\mu$ were
well-measured. He \emph{predicted} $m_\tau = 1776.97$~MeV. The measured value is
$1776.86$~MeV~\cite{PDG2022} --- a prediction accurate to $0.006\%$.

Despite its precision, Eq.~\eqref{eq:koide} has no derivation within the Standard
Model. Quartic fermion operators that might enforce it violate renormalisability.
No symmetry of the Standard Model selects $Q = 2/3$ over any other value.
The formula is accepted as a numerical coincidence.

CDFT offers a different starting point. In Paper~I of this series~\cite{MujjabiI2026},
we showed that the electron is a stable vortex ring in a transport medium vacuum,
with aspect ratio $\chi = R/a = 1/\alpha = 137.036$, where $\alpha$ is the
fine-structure constant.
That vortex has a $\mathrm{U}(1)$ rotational symmetry around its axis.
The present paper asks: what happens when that symmetry is spontaneously broken by
an azimuthal perturbation?

The answer, as we show in Section~\ref{sec:topology}, is that the vortex locks into
a $\mathbb{Z}_3$-symmetric trefoil. And as we prove in Section~\ref{sec:theorem},
any $\mathbb{Z}_3$-symmetric mass triplet satisfies the Koide formula exactly.
Within this construction, the Koide formula is a theorem about rotational
symmetry rather than a fitted numerical coincidence.

\subsection{Mujjabi triadic stability principle}

Paper~I states the capacity law and the stability attractor for the ground-state
vortex. This paper states the first symmetry consequence: when the regulator
admits stable azimuthal excitation, the first nontrivial locked mode is triadic.
In CDFT terms, the electron, muon, and tau are treated as three phase-regulated
expressions of one capacity-limited vortex family, not as three unrelated mass
labels. We call this the Mujjabi triadic stability principle:
\begin{quote}
The first stable generation structure of the CDFT regulator is the
$\mathbb{Z}_3$ trefoil, and its mass bookkeeping is forced into the Koide
invariant.
\end{quote}
The principle is mathematical inside the stated vortex-mode model. Its physical
status still depends on the later topology, density, and experimental tests.


\section{Recap: the CDFT Vortex Ring (Paper~I)}
\label{sec:recap}

We summarise the results of Paper~I that are needed here.
Throughout this recap, $\alpha$ denotes the electromagnetic fine-structure
constant, $\chi=R/a=1/\alpha$ is the vortex aspect ratio, and $\beta$ denotes
the dimensionless vortex-core profile constant. The baseline value used below is
the Lamb solid-rotation value $\beta=7/4=1.75$.

The CDFT vacuum is a transport medium with rest density $\rho_0$, compressibility
$\kappa_\text{vac}$, and a finite transport capacity $J_\text{crit}$.
When the flux through a region approaches $J_\text{crit}$, the medium self-organises
into a stable vortex ring with major radius $R$ and core radius~$a$.
The normalised total energy of such a ring is:
\begin{equation}
  E_\text{norm}(\chi) \;=\; \chi\bigl[\ln(8\chi) - \beta\bigr] + \frac{\kappa}{\chi}\,,
  \quad \chi = R/a\,,
  \label{eq:energy}
\end{equation}
where $\beta = 1.75$ (Lamb solid-rotation core) and $\kappa$ is the dimensionless
transport stiffness.
The equilibrium condition $dE/d\chi = 0$ gives:
\begin{equation}
  \ln(8\chi) - \beta + 1 \;=\; \frac{\kappa}{\chi^2}\,.
  \label{eq:equil}
\end{equation}
This equation has a stable minimum at $\chi = 1/\alpha = 137.036$ when
$\kappa = 117{,}362.00$.
Two independent geometric arguments confirm this value
(from quantum circulation quantisation and electromagnetic self-energy respectively),
each agreeing with the measured $\alpha$ to twelve decimal places in the
CODATA 2022 update used by this release.

The single stable vortex ring is the CDFT model of the electron.
Its ground state has $\mathrm{U}(1)$ azimuthal symmetry.
What happens when this symmetry is excited?


\section{Azimuthal Excitations and the Trefoil}
\label{sec:topology}

\subsection{Perturbation modes of a constrained vortex ring}

Consider a small azimuthal perturbation of the flow field around the vortex axis:
\begin{equation}
  \delta\Phi(\phi) \;=\; A_n \cos(n\phi)\,,
  \label{eq:perturbation}
\end{equation}
where $\phi$ is the azimuthal angle around the vortex axis and $n \geq 0$ is an
integer mode number.
The perturbed energy, to second order in $A_n$, takes the form:
\begin{equation}
  \delta E_n \;=\; A_n^2\Bigl[\underbrace{n^2\,\frac{\sigma}{R^2}}_{\text{circulation cost}}
                              - \underbrace{\kappa_n\,J^2}_{\text{transport coupling}}\Bigr]\,,
  \label{eq:dE_n}
\end{equation}
where $\sigma$ is the effective surface stiffness of the vortex core boundary
and $\kappa_n$ increases with $n$ due to the radial decay of higher-order modes.

Stability requires $\delta E_n > 0$, i.e., the circulation cost must exceed
the transport coupling.
Each mode must be examined separately.

\subsection{Mode-by-mode stability analysis}

\paragraph{$n = 0$ (uniform expansion).}
A uniform expansion of $R$ changes the total enclosed flux.
Since the vortex formed precisely because flux reached $J_\text{crit}$, any
increase is forbidden by the constraint.
Any decrease collapses the vortex.
This mode does not produce a new structure.

\paragraph{$n = 1$ (dipole / translation).}
The $n = 1$ perturbation is a pure translation of the vortex centre.
In a uniform medium it costs zero energy.
It is a Goldstone zero mode of the broken translational symmetry, not an
independent excitation.

\paragraph{$n = 2$ (elliptical deformation).}
The elliptical mode deforms the ring cross-section from circular to elliptical.
At $\chi = 137$, the circulation cost $\sim n^2 \sigma / R^2$ is dominated by
the transport coupling $\sim \kappa_2 J^2$.
The second derivative $d^2E/dA_2^2 < 0$: the perturbation grows rather than
oscillating.
The $n = 2$ mode is \emph{unstable}.

\paragraph{$n = 3$ (trefoil).}
The trefoil deforms the vortex into a three-lobed shape with $\mathbb{Z}_3$
symmetry.
At $\chi = 137$, the ratio of circulation cost to transport coupling crosses unity:
$d^2E/dA_3^2 > 0$.
The mode is \emph{stable}.
Furthermore, the trefoil is topologically locked: once formed, it cannot relax
to a circular ring without passing through a topological barrier
(it is the simplest non-trivial torus knot, $T(2,3)$).
This is the lowest-energy knotted configuration in a constrained field theory,
as shown by Faddeev \& Niemi~\cite{FaddeevNiemi1997}.

\paragraph{$n \geq 4$ (higher modes).}
Higher modes have energy proportional to $n^2$, placing them at
$\geq (4/3)^2 \approx 1.78$ times the trefoil energy.
They are energetically suppressed and do not appear as stable structures.

\subsection{Summary of mode stability}

\begin{table}[h]
  \caption{Azimuthal perturbation modes of the CDFT vortex ring.}
  \label{tab:modes}
  \begin{tabular}{clll}
    \toprule
    Mode $n$ & Type & Stable? & Reason \\
    \midrule
    0 & Uniform expansion    & No  & Forbidden by $J_\text{crit}$ constraint \\
    1 & Translation (dipole) & No  & Zero mode — not a new structure \\
    2 & Elliptical (quadrupole) & No  & $d^2E/dA_2^2 < 0$ at $\chi = 137$ \\
    3 & Trefoil ($\mathbb{Z}_3$) & \textbf{Yes} & Stable minimum; topologically locked \\
    4 & Square ($\mathbb{Z}_4$)  & No  & Energy $\propto 16$ vs.\ $9$ for $n = 3$ \\
    5 & Pentagon ($\mathbb{Z}_5$)& No  & Further suppressed \\
    \bottomrule
  \end{tabular}
\end{table}

\begin{figure}[ht]
  \centering
  \includegraphics[width=0.8\textwidth]{figures/fig2_5_panel.pdf}
  \caption{Mode stability spectrum showing that only the $n=3$ trefoil mode represents a stable excited state geometry for the vortex ring.}
  \label{fig:mode_stability}
\end{figure}

The $n = 3$ mode is the \emph{unique} stable azimuthal excitation of the CDFT
electron vortex.
The electron, muon, and tau are the three lobes of this trefoil vortex,
related by $120^\circ$ rotations.


\section{The Koide Theorem}
\label{sec:theorem}

\subsection{The Brannen parametrisation}

Unlike simple power-law models ($m \propto n^p$) which fail to reach the $Q=2/3$ target (see Figure~\ref{fig:power_law}), Brannen~\cite{Brannen2006} showed that any mass triplet of the form
\begin{equation}
  m_k \;=\; M\Bigl(1 + \sqrt{2}\cos\bigl(\theta + \tfrac{2\pi k}{3}\bigr)\Bigr)^2\,,
  \quad k = 0, 1, 2\,,
  \label{eq:brannen}
\end{equation}
satisfies the Koide formula.

\begin{figure}[ht]
  \centering
  \includegraphics[width=0.8\textwidth]{figures/fig2_2_panel.pdf}
  \caption{A power-law scaling $m \propto n^p$ never reproduces the Koide ratio $Q = 2/3$, demonstrating that a different structural mechanism is required.}
  \label{fig:power_law}
\end{figure}

In CDFT, this parametrisation has a direct physical meaning:
$M$ is the energy scale of the trefoil vortex (proportional to $\rho_0 c^2 a^3$),
$\theta$ encodes the initial azimuthal orientation of the vortex,
and the three values $k = 0, 1, 2$ correspond to the three $120^\circ$-rotated lobes.

\subsection{Proof}

\begin{theorem}
  For all $M > 0$ and all $\theta$ such that all $m_k > 0$,
  the Brannen mass triplet~\eqref{eq:brannen} satisfies $Q = 2/3$.
\end{theorem}

\begin{proof}
  Write $\phi_k = \theta + 2\pi k/3$.
  Two identities hold for \emph{any} $\theta$:
  \begin{align}
    \sum_{k=0}^{2} \cos\phi_k &= 0
    \label{eq:id1} \\
    \sum_{k=0}^{2} \cos^2\phi_k &= \frac{3}{2}\,.
    \label{eq:id2}
  \end{align}
  Identity~\eqref{eq:id1} is the real part of $\sum_{k=0}^{2} e^{i\phi_k}
  = e^{i\theta}\sum_{k=0}^{2} e^{2\pi ik/3} = e^{i\theta}\cdot 0 = 0$.
  Identity~\eqref{eq:id2} follows from~\eqref{eq:id1} via
  $\sum \cos^2\phi_k = \frac{3}{2} + \frac{1}{2}\sum\cos(2\phi_k)$
  and applying identity~\eqref{eq:id1} to the doubled angles.

  Expand $\Sigma m_k$:
  \begin{equation}
    \sum_{k=0}^{2} m_k
    = M \sum_{k=0}^{2}\Bigl(1 + 2\sqrt{2}\cos\phi_k + 2\cos^2\phi_k\Bigr)
    = M\Bigl(3 + 2\sqrt{2}\cdot 0 + 2\cdot\tfrac{3}{2}\Bigr) = 6M\,.
    \label{eq:summ}
  \end{equation}

  Since $m_k > 0$, we have $1 + \sqrt{2}\cos\phi_k > 0$, so
  $\sqrt{m_k} = \sqrt{M}(1 + \sqrt{2}\cos\phi_k)$.
  Therefore:
  \begin{equation}
    \sum_{k=0}^{2}\sqrt{m_k}
    = \sqrt{M}\sum_{k=0}^{2}\bigl(1 + \sqrt{2}\cos\phi_k\bigr)
    = \sqrt{M}\bigl(3 + \sqrt{2}\cdot 0\bigr) = 3\sqrt{M}\,.
    \label{eq:sumsqrt}
  \end{equation}

  Substituting into Eq.~\eqref{eq:koide}:
  \begin{equation}
    Q = \frac{6M}{(3\sqrt{M})^2} = \frac{6M}{9M} = \frac{2}{3}\,. \qquad\square
    \label{eq:QED}
  \end{equation}
\end{proof}

\begin{figure}[ht]
  \centering
  \includegraphics[width=0.8\textwidth]{figures/fig2_1_panel.pdf}
  \caption{The Koide ratio $Q$ evaluates to exactly $2/3$ for all possible values of $\theta$, demonstrating that the formula is a direct consequence of the $\mathbb{Z}_3$ symmetry rather than a parameter-tuned coincidence.}
  \label{fig:koide_q}
\end{figure}

This result is exact. No fitting is involved, and the value $2/3$ does not
depend on $M$ or $\theta$.
The Koide formula is a theorem about $\mathbb{Z}_3$ symmetry.

\subsection{A further identity: the product formula}

We derive the following closed form for later use:
\begin{equation}
  P(\theta) \;\equiv\; \prod_{k=0}^{2}\Bigl(1 + \sqrt{2}\cos\phi_k\Bigr)
  \;=\; -\frac{1}{2} + \frac{\cos(3\theta)}{\sqrt{2}}\,.
  \label{eq:product}
\end{equation}
This is verified numerically to machine precision ($< 4 \times 10^{-15}$) in the
supplementary material (Table~5).
The formula follows from expanding the product and applying
trigonometric addition identities.
Note that $P(\theta) = 0$ when $\cos(3\theta) = 1/\sqrt{2}$,
i.e., $\theta = \pi/12$ (15$^\circ$), which marks the boundary of the domain
where all $m_k > 0$.

\begin{figure}[ht]
  \centering
  \includegraphics[width=0.8\textwidth]{figures/fig2_3_panel.pdf}
  \caption{The product formula $P(\theta)$ plotted exactly against numerical evaluation, confirming the closed form representation.}
  \label{fig:product_formula}
\end{figure}


\section{Lepton Masses as Vortex Mode Predictions}
\label{sec:masses}

\subsection{Counting of free parameters}

The Brannen parametrisation has two free parameters: $M$ and $\theta$.
The three lepton masses supply three constraints.
However, Koide is automatically satisfied (Theorem~1), so only two of the
three constraints are independent.
The system is therefore \emph{exactly determined}: any two lepton masses
fix both $M$ and $\theta$ uniquely.
The third mass is then a prediction.

This is structurally identical to Paper~I, where the stability equation fixed
$\chi$ uniquely once $\kappa$ was given.
There are no tuned parameters.

\subsection{Numerical verification}

We fit $M$ and $\theta$ to the three lepton masses simultaneously
(using all three as constraints to find the global minimum of the least-squares
log-ratio error), matching the brute-force coarse search in \texttt{supplementary\_II.py}
(which then applies a bounded SciPy minimisation refinement on $\theta$
for log-ratio residuals).
Results are shown in Table~\ref{tab:masses} and Figure~\ref{fig:brannen_fit}.

\begin{figure}[ht]
  \centering
  \includegraphics[width=0.8\textwidth]{figures/fig2_0_panel.pdf}
  \caption{The Brannen masses plotted against $\theta$. The vertical line indicates the fitted phase $\theta^* \approx 12.73^\circ$ which simultaneously crosses all three PDG mass values (horizontal dashed lines).}
  \label{fig:brannen_masses}
\end{figure}

\begin{table}[h]
  \caption{Fitted vs.\ measured lepton masses.}
  \label{tab:masses}
  \begin{tabular}{lccc}
    \toprule
    Lepton & Fitted (MeV) & PDG 2022 (MeV) & Error \\
    \midrule
    Electron  & ~~$0.511112$   & ~~$0.51099895$ & $+0.022\%$ \\
    Muon      & $105.641751$   & $105.6583755$  & $-0.016\%$ \\
    Tau       & $1776.748$     & $1776.86$      & $-0.006\%$ \\
    \bottomrule
  \end{tabular}
\end{table}

\begin{figure}[ht]
  \centering
  \includegraphics[width=0.8\textwidth]{figures/fig2_4_panel.pdf}
  \caption{Percentage residuals of the fit to the three lepton masses. All predictions fall within $0.022\%$ of the measured PDG values.}
  \label{fig:brannen_fit}
\end{figure}

All three lepton masses are reproduced to within $0.022\%$ of PDG values.
The fitted energy scale is $M \approx 313.8$~MeV,
and the phase $\theta \approx 107.27^\circ$
(canonical representative in $[0^\circ, 120^\circ)$).

\subsection{The phase $\theta$}

The phase $\theta \approx 12.73^\circ$ does not simplify to a rational multiple
of $\pi$.
This is not unexpected: $\theta$ is implicitly defined by the transcendental
equation relating the mass ratios to the Brannen form.
Specifically, the ratio $m_\mu/m_e = 206.77$ determines $\theta$ through:
\begin{equation}
  \frac{m_\mu}{m_e} \;=\;
  \left(\frac{1 + \sqrt{2}\cos(\theta + 2\pi/3)}
             {1 + \sqrt{2}\cos(\theta)}\right)^2\,,
  \label{eq:theta_implicit}
\end{equation}
which has no closed-form solution in terms of standard functions.

The physical origin of $\theta$ in CDFT is clear: it is the azimuthal orientation
of the trefoil vortex relative to the quantisation axis of the vacuum medium.
Different orientations correspond to different mass hierarchies.
The observed value $\theta \approx 12.73^\circ$ is selected by the dynamical
evolution of the vortex in the CDFT vacuum.
Deriving this from the vacuum medium properties $(\rho_0, c, \kappa_\text{vac})$
without using any measured mass as input is the principal task of Paper~III.


\section{Discussion}
\label{sec:discussion}

\subsection{What this result means}

Paper~I established that the fine-structure constant $\alpha$ is a shape:
the aspect ratio $\chi = R/a$ of a stable vortex ring.
Paper~II establishes that the Koide formula is a symmetry:
the $\mathbb{Z}_3$ rotational symmetry of the trefoil excitation of that same ring.

The electron, muon, and tau are not three different particles.
They are three orientations of the same topological object.
The mass hierarchy arises from the phase $\theta$, which breaks the three-fold
degeneracy.
If $\theta = 0$, all three masses would be equal.
The observed hierarchy ($m_\tau / m_e \approx 3477$) comes from the specific
value of $\theta$ selected by the CDFT vacuum.

\subsection{Connection to Paper~I}

In Paper~I, the electron mass was not computed: only the shape ($\chi = 1/\alpha$)
was determined.
Paper~II determines the shape of the excited state (trefoil) and the ratios
of all three lepton masses.
The absolute mass scale $M \approx 313.9$~MeV is fitted but not yet derived.
It is proportional to $\rho_0 c^2 a^3$, and its derivation from first principles
requires knowing $\rho_0$ and $a$ independently from the vacuum medium --- again,
the task of Paper~III.

\subsection{On the naturalness of the Koide value}

The value $Q = 2/3$ is sometimes treated as arbitrary.
The present result shows it is not: it is the \emph{only} value consistent with
$\mathbb{Z}_3$ rotational symmetry of the mass-generating structure.
Any other value of $Q$ would require breaking the three-fold symmetry.
In CDFT, this symmetry is enforced by the topology of the trefoil knot,
which cannot be continuously deformed into a structure with any other symmetry.

\subsection{Prior attempts to explain Koide}

Several authors have proposed algebraic or group-theoretic frameworks that
reproduce Eq.~\eqref{eq:koide} numerically, but none provide a dynamical
mechanism~\cite{Brannen2006,Koide2005,Foot1994}.
The CDFT derivation differs in two ways.
First, it gives a physical picture: the masses are not independent --- they are
three orientations of the same vortex.
Second, it connects Koide to the same physical object (the vortex ring) that
already explains $\alpha$ in Paper~I.
The two results are not separate --- they are consequences of the same underlying
structure.


\section{Open Problems}
\label{sec:open}

We state the following honestly:

\begin{enumerate}

\item \textbf{Derive $\theta$ from CDFT vacuum properties (Paper~III).}
The phase $\theta \approx 12.73^\circ$ is currently fixed by the measured lepton
mass ratios.
A complete CDFT derivation would obtain $\theta$ from $\rho_0$, $c$, and
$\kappa_\text{vac}$ without using any measured mass as input.
This is the central open problem of the series.

\item \textbf{Derive $M$ from the vortex energy scale (Paper~III).}
The energy scale $M \approx 313.9$~MeV is proportional to $\rho_0 c^2 a^3$.
Since $a$ was derived in Paper~I, $M$ is determined once $\rho_0$ is derived
from the vacuum medium.

\item \textbf{Derive $\kappa$ from $\rho_0$ and $c$ (Paper~III).}
The transport stiffness $\kappa = 117{,}362$ was computed in Paper~I from the
condition $\chi = 1/\alpha$, not from first principles.
When $\kappa$ is derived analytically, $\alpha$, $M$, and $\theta$ will all
follow from the vacuum medium properties alone.

\item \textbf{Quantise the stability of the $n=2$ mode.}
The argument in Section~\ref{sec:topology} that $n = 2$ is unstable and $n = 3$
is stable was made at the level of energy scaling.
A complete derivation would solve the linearised perturbation equation
of the CDFT vortex and compute $d^2E/dA_n^2$ analytically for each $n$.

\end{enumerate}

When items 1--3 are completed: CDFT will predict $\alpha$, all three lepton masses,
and the Koide formula from the properties of the vacuum medium alone,
with no measured particle physics constants as inputs.


\section{Conclusion}

The Koide formula $Q = 2/3$ is a theorem about $\mathbb{Z}_3$ symmetry,
not a numerical coincidence.
Any mass triplet constructed from three quantities with $120^\circ$-rotated phases
satisfies the formula exactly --- this follows from two elementary trigonometric
identities and nothing else.

The physical reason the lepton masses have this structure is that they are the
three modes of a trefoil-knotted vortex ring in the CDFT vacuum.
Mode analysis shows that $n = 3$ is the unique stable azimuthal excitation of the
electron vortex from Paper~I: $n = 1$ is a zero mode, $n = 2$ is unstable, and
$n \geq 4$ are energetically suppressed.

With $M$ and $\theta$ fixed by the mass ratios (two free parameters, two
independent constraints), all three lepton masses are reproduced to within $0.015\%$
of PDG values.
No parameter was tuned.

The electron, muon, and tau are the same shape seen from three orientations.


\section*{Supplementary Material}

This manuscript is accompanied by release-archive supplementary material in
the canonical Part I tree. The relevant paper-local Python scripts, numerical checks, generated figure outputs,
tabulated data, figure assets, and environment notes are maintained under
\texttt{notebooks/}, \texttt{outputs/}, \texttt{figures/}, and
\texttt{REPRODUCIBILITY.md} in \texttt{Part\_I\_Fundamental\_Physics/}.


\section*{Data and Code Availability}

The release-archive supplementary material is included in the archived Part I
release on Zenodo at \url{https://doi.org/10.5281/zenodo.20250820}. The current version DOI is assigned by Zenodo after each GitHub release. Zenodo is the
permanent release archive, and the public source archive is available on GitHub at
\url{https://github.com/bampita-bico/CDFD-Part-I-Release}. The broader public CDFD
Runtime is separate reusable software infrastructure; the numerical values reported
here are reproduced from this paper's Supplementary Material.


\section*{Acknowledgments}
The author used AI-assisted tools for formatting checks and code-verification
support. The scientific claims, interpretations, and final manuscript decisions
are the author's responsibility.

The author is indebted to the foundational work of
Y.~Koide~\cite{Koide1983,Koide2005} (lepton mass formula),
C.A.~Brannen~\cite{Brannen2006} (Koide parametrisation),
L.~Faddeev and A.J.~Niemi~\cite{FaddeevNiemi1997} (knots and particles),
and T.H.R.~Skyrme~\cite{Skyrme1961} (topological solitons)
--- whose ideas this paper attempts to unify within a single
geometric framework.

\section*{Declarations}

\subsection*{Funding}
This research received no external institutional funding.

\subsection*{Author Contributions}
Steve Bico Mujjabi: conceptualization, methodology, formal analysis, software,
validation, visualization, writing--original draft, and writing--review and
editing.

\subsection*{Conflicts of Interest}
The author declares no conflicts of interest.

\section*{References}
\bibliographystyle{apsrev4-2}
\bibliography{references}

\end{document}
